\begin{thebibliography}{99}

\bibitem{bai13} R. A. Bailey. \emph{Latin squares}. G. C. Steward lecture, Gonville and Caius College, Cambridge, 7 March 2013 (slides).

\bibitem{dey17} A. Dey. \emph{The problem of 36 officers, Euler's conjecture and related matters}. Special Proceedings of the 19th Annual Conference of SSCA (2017), 1--7.

\bibitem{oza1723} J. Ozanam. \emph{R\'ecr\'eations math\'ematiques et physiques}, vol.~1, 1723.

\bibitem{eul1782} L. Euler. \emph{Recherches sur une nouvelle esp\`ece de quarr\'es magiques}. Verh. Zeeuwsch. Genootsch. Wetensch. Vlissingen 9 (1782), 85--239.

\bibitem{tar1900} G. Tarry. \emph{Le probl\`eme des 36 officiers}. Secr\'etariat de l'Association fran\c{c}aise pour l'avancement des sciences, 1900.

\bibitem{bs59} R. C. Bose, S. S. Shrikhande. \emph{On the falsity of Euler's conjecture about the non-existence of two orthogonal Latin squares of order $4t+2$}. Proc. Natl. Acad. Sci. USA 45 (1959), 734--737.

\bibitem{bsp60} R. C. Bose, S. S. Shrikhande, E. T. Parker. \emph{Further results on the construction of mutually orthogonal Latin squares and the falsity of Euler's conjecture}. Canad. J. Math. 12 (1960), 189--203.

\bibitem{wan07} I. M. Wanless. \emph{Transversals in Latin squares}. Quasigroups Related Systems 15 (2007), 169--190. \href{https://arxiv.org/abs/0903.5142}{arXiv:0903.5142}

\bibitem{rys67} H. J. Ryser. \emph{Neuere Probleme der Kombinatorik}. Vortr\"age \"uber Kombinatorik, Oberwolfach (1967), 69--91.

\bibitem{ste75} S. K. Stein. \emph{Transversals of Latin squares and their generalizations}. Pacific J. Math. 59 (1975), 567--575.

\bibitem{mon23} R. Montgomery. \emph{A proof of the Ryser--Brualdi--Stein conjecture for large even $n$}. \href{https://arxiv.org/abs/2310.19779}{arXiv:2310.19779}

\bibitem{fy} R. A. Fisher, F. Yates. \emph{Statistical Tables for Agricultural, Biological, and Medical Research}. Oliver and Boyd, Edinburgh.

\bibitem{bose38} R. C. Bose. \emph{On the application of the properties of Galois fields to the problem of construction of hyper-Graeco-Latin squares}. Sankhy\=a (1938).

\bibitem{ste39} W. L. Stevens. \emph{The completely orthogonalized Latin square}. Annals of Eugenics (1939).

\bibitem{man42} H. B. Mann. \emph{The construction of orthogonal Latin squares}. Ann. Math. Statistics 13 (1942), 418--423.

\bibitem{man44} H. B. Mann. \emph{On orthogonal Latin squares}. Bull. Amer. Math. Soc. 50 (1944), 249--257.

\bibitem{pai47} L. J. Paige. \emph{A note on finite abelian groups}. Bull. Amer. Math. Soc. 53 (1947), 590--593.

\bibitem{bat50} P. Bateman. \emph{Complete mappings of infinite groups}. Amer. Math. Monthly 57 (1950), 621--622.

\bibitem{pai51} L. J. Paige. \emph{Complete mappings of finite groups}. Pacific J. Math. 1 (1951), 111--116.

\bibitem{mon24} R. Montgomery. \emph{Transversals in Latin squares}. In: Surveys in Combinatorics 2024, London Math. Soc. Lecture Note Ser., Cambridge Univ. Press, 131--158. \href{https://arxiv.org/abs/2406.19873}{arXiv:2406.19873}

\bibitem{hall52} M. Hall. \emph{A combinatorial problem on abelian groups}. Proc. Amer. Math. Soc. 3 (1952), 584--587.

\bibitem{hp55} M. Hall, L. J. Paige. \emph{Complete mappings of finite groups}. Pacific J. Math. 5 (1955), 541--549.

\bibitem{wil09} S. Wilcox. \emph{Reduction of the Hall--Paige conjecture to sporadic simple groups}. J. Algebra 321 (2009), 1407--1428. \href{https://arxiv.org/abs/1010.1323}{arXiv:1010.1323}

\bibitem{eva09} A. B. Evans. \emph{The admissibility of sporadic simple groups}. J. Algebra 321 (2009), 105--116.

\bibitem{bccsz20} J. N. Bray, Q. Cai, P. J. Cameron, P. Spiga, H. Zhang. \emph{The Hall--Paige conjecture, and synchronization for affine and diagonal groups}. J. Algebra 545 (2020), 27--42. \href{https://arxiv.org/abs/1811.12671}{arXiv:1811.12671}

\bibitem{emm22} S. Eberhard, F. Manners, R. Mrazovi\'c. \emph{An asymptotic for the Hall--Paige conjecture}. Adv. Math. 404 (2022), 108423. \href{https://arxiv.org/abs/2003.01798}{arXiv:2003.01798}

\bibitem{mp25a} A. M\"uyesser, A. Pokrovskiy. \emph{A random Hall--Paige conjecture}. Invent. Math. 240 (2025), 779--867. \href{https://arxiv.org/abs/2204.09666}{arXiv:2204.09666}

\bibitem{gh20} L. Goddyn, K. Halasz. \emph{All group-based Latin squares possess near transversals}. J. Combin. Des. 28 (2020), 358--365. \href{https://arxiv.org/abs/1905.07071}{arXiv:1905.07071}

\bibitem{sne99} H. S. Snevily. \emph{The Cayley addition table of $\mathbb Z_n$}. Amer. Math. Monthly 106 (1999), 584--585.

\bibitem{alon00} N. Alon. \emph{Additive Latin transversals}. Israel J. Math. 117 (2000), 125--130.

\bibitem{dkss01} S. Dasgupta, G. K\'arolyi, O. Serra, B. Szegedy. \emph{Transversals of additive Latin squares}. Israel J. Math. 126 (2001), 17--28.

\bibitem{ars11} B. Arsovski. \emph{A proof of Snevily's conjecture}. Israel J. Math. 182 (2011), 505--508.

\bibitem{aa04} S. Akbari, A. Alipour. \emph{Transversals and multicolored matchings}. J. Combin. Des. 12 (2004), 325--332.

\bibitem{var91} I. Vardi. \emph{Computational Recreations in Mathematica}. Addison-Wesley, 1991.

\bibitem{emm19} S. Eberhard, F. Manners, R. Mrazovi\'c. \emph{Additive triples of bijections, or the toroidal semiqueens problem}. J. Eur. Math. Soc. (2019), 441--463. \href{https://arxiv.org/abs/1510.05987}{arXiv:1510.05987}

\bibitem{pol1918} G. P\'olya. \emph{\"Uber die ``doppelt-periodischen'' L\"osungen des $n$-Damen-Problems}. In: W. Ahrens (ed.), Mathematische Unterhaltungen und Spiele 2, Teubner, 1918, 364--374.

\bibitem{bez1848} M. Bezzel (``Schachfreund''). Proposal of the eight-queens problem. Berliner Schachzeitung 3 (1848), 363.

\bibitem{nau1850} F. Nauck. \emph{Briefwechseln mit allen f\"ur alle}. Illustrirte Zeitung 15 (1850), 182.

\bibitem{pau1874} E. Pauls. \emph{Das Maximalproblem der Damen auf dem Schachbrete}. Deutsche Schachzeitung 29 (1874), 129--134, 257--267.

\bibitem{rvz94} I. Rivin, I. Vardi, P. Zimmermann. \emph{The $n$-queens problem}. Amer. Math. Monthly 101 (1994), 629--639.

\bibitem{lur17} Z. Luria. \emph{New bounds on the number of $n$-queens configurations}. \href{https://arxiv.org/abs/1705.05225}{arXiv:1705.05225}

\bibitem{ls21} Z. Luria, M. Simkin. \emph{A lower bound for the $n$-queens problem}. \href{https://arxiv.org/abs/2105.11431}{arXiv:2105.11431}

\bibitem{sim21} M. Simkin. \emph{The number of $n$-queens configurations}. Adv. Math. (2023). \href{https://arxiv.org/abs/2107.13460}{arXiv:2107.13460}

\bibitem{gms22} S. Glock, D. Munh\'a Correia, B. Sudakov. \emph{The $n$-queens completion problem}. \href{https://arxiv.org/abs/2111.11402}{arXiv:2111.11402}

\bibitem{dk24} T. Dai, T. Kelly. \emph{On the existence of reflecting $n$-queens configurations}. Forum Math. Sigma (2024). \href{https://arxiv.org/abs/2407.12742}{arXiv:2407.12742}

\bibitem{eva13} A. B. Evans. \emph{The existence of strong complete mappings of finite groups: a survey}. Discrete Math. 313 (2013), 1191--1196.

\bibitem{fgt81} R. J. Friedlander, B. Gordon, P. Tannenbaum. \emph{Partitions of groups and complete mappings}. Pacific J. Math. 92 (1981), 283--293.

\bibitem{muy23} A. M\"uyesser. \emph{Cycle type in Hall--Paige: a proof of the Friedlander--Gordon--Tannenbaum conjecture}. Forum Math. Sigma 14 (2026), e50. \href{https://arxiv.org/abs/2303.16157}{arXiv:2303.16157}

\bibitem{ks02} A. E. K\'ezdy, H. S. Snevily. \emph{Distinct sums modulo $n$ and tree embeddings}. Combin. Probab. Comput. 11 (2002), 35--42.

\bibitem{bk21} C. Bowtell, P. Keevash. \emph{The $n$-queens problem}. \href{https://arxiv.org/abs/2109.08083}{arXiv:2109.08083}

\end{thebibliography}
